% GENERATED FILE. Edit canonical lessons, then run npm run generate:books.
% canonical-source-hash: fnv1a64:a96df1e95b96121f
% canonical-lessons: TE-C238-jnapakam, TE-C238-caritra, TE-C238-bhugolasastram, TE-C238-vijnanasastram, TE-C238-bhautikasastram

\chapter{Memory, History, Geography, Science, Physics}
\label{ch:te-memory-history-geography-science-physics}
\hlchaptermodality{238}

\begin{chapteropening}
\textbf{By the end of this chapter:} \emph{I can say a memory, history, geography, science and physics.}

Complete the last lesson of chapter 238: I can say a memory, history, geography, science and physics.
\end{chapteropening}

\section[\texorpdfstring{\te{జ్ఞాపకం} (jñāpakaṁ)}{జ్ఞాపకం (jñāpakaṁ)}]{\te{జ్ఞాపకం} (jñāpakaṁ) — a memory}

\label{lesson:TE-C238-jnapakam}



\begin{quote}
Before the new one: say the Telugu for a bus conductor, then the Telugu for a lawyer.
\end{quote}

\subsection*{You'll want to know: \te{జ్ఞాపకం}}
\textbf{\te{జ్ఞాపకం}} — \emph{jñāpakaṁ} — \textquotedblleft{}a memory\textquotedblright{}.

\begin{grammarlens}[title={What you've built}]
One more word for this chapter.
\end{grammarlens}

\subsection*{Your turn}
\begin{itemize}
\raggedright
  \item \emph{Say it:} \emph{jñāpakaṁ}
  \item \emph{Say it:} \emph{jñāpakaṁ}, once more
  \item \emph{Say it:} say \emph{jñāpakaṁ}
  \item \emph{Recall:} say the Telugu for a bus conductor, then the Telugu for a lawyer, then say \emph{jñāpakaṁ} again
\end{itemize}

\begin{tcolorbox}[breakable,colback=teal!4,colframe=teal!35!black,title={Before you move on}]
What does \te{జ్ఞాపకం} mean? (\textquotedblleft{}A memory\textquotedblright{}.) Say it once more.
\end{tcolorbox}
\section[\texorpdfstring{\te{చరిత్ర} (caritra)}{చరిత్ర (caritra)}]{\te{చరిత్ర} (caritra) — history}

\label{lesson:TE-C238-caritra}



\begin{quote}
Before the new one: say the Telugu for a lawyer, then the Telugu for a memory.
\end{quote}

\subsection*{You'll want to know: \te{చరిత్ర}}
\textbf{\te{చరిత్ర}} — \emph{caritra} — \textquotedblleft{}history\textquotedblright{}.

\begin{grammarlens}[title={What you've built}]
One more word for this chapter.
\end{grammarlens}

\subsection*{Your turn}
\begin{itemize}
\raggedright
  \item \emph{Say it:} \emph{caritra}
  \item \emph{Say it:} \emph{caritra}, once more
  \item \emph{Say it:} say \emph{caritra}
  \item \emph{Recall:} say the Telugu for a lawyer, then the Telugu for a memory, then say \emph{caritra} again
\end{itemize}

\begin{tcolorbox}[breakable,colback=teal!4,colframe=teal!35!black,title={Before you move on}]
What does \te{చరిత్ర} mean? (\textquotedblleft{}History\textquotedblright{}.) Say it once more.
\end{tcolorbox}
\section[\texorpdfstring{\te{భూగోళశాస్త్రం} (bhūgōḷaśāstraṁ)}{భూగోళశాస్త్రం (bhūgōḷaśāstraṁ)}]{\te{భూగోళశాస్త్రం} (bhūgōḷaśāstraṁ) — geography}

\label{lesson:TE-C238-bhugolasastram}



\begin{quote}
Before the new one: say the Telugu for a memory, then the Telugu for history.
\end{quote}

\subsection*{You'll want to know: \te{భూగోళశాస్త్రం}}
\textbf{\te{భూగోళశాస్త్రం}} — \emph{bhūgōḷaśāstraṁ} — \textquotedblleft{}geography\textquotedblright{}.

\begin{grammarlens}[title={What you've built}]
One more word for this chapter.
\end{grammarlens}

\subsection*{Your turn}
\begin{itemize}
\raggedright
  \item \emph{Say it:} \emph{bhūgōḷaśāstraṁ}
  \item \emph{Say it:} \emph{bhūgōḷaśāstraṁ}, once more
  \item \emph{Say it:} say \emph{bhūgōḷaśāstraṁ}
  \item \emph{Recall:} say the Telugu for a memory, then the Telugu for history, then say \emph{bhūgōḷaśāstraṁ} again
\end{itemize}

\begin{tcolorbox}[breakable,colback=teal!4,colframe=teal!35!black,title={Before you move on}]
What does \te{భూగోళశాస్త్రం} mean? (\textquotedblleft{}Geography\textquotedblright{}.) Say it once more.
\end{tcolorbox}
\section[\texorpdfstring{\te{విజ్ఞానశాస్త్రం} (vijñānaśāstraṁ)}{విజ్ఞానశాస్త్రం (vijñānaśāstraṁ)}]{\te{విజ్ఞానశాస్త్రం} (vijñānaśāstraṁ) — science}

\label{lesson:TE-C238-vijnanasastram}



\begin{quote}
Before the new one: say the Telugu for history, then the Telugu for geography.
\end{quote}

\subsection*{You'll want to know: \te{విజ్ఞానశాస్త్రం}}
\textbf{\te{విజ్ఞానశాస్త్రం}} — \emph{vijñānaśāstraṁ} — \textquotedblleft{}science\textquotedblright{}.

\begin{grammarlens}[title={What you've built}]
One more word for this chapter.
\end{grammarlens}

\subsection*{Your turn}
\begin{itemize}
\raggedright
  \item \emph{Say it:} \emph{vijñānaśāstraṁ}
  \item \emph{Say it:} \emph{vijñānaśāstraṁ}, once more
  \item \emph{Say it:} say \emph{vijñānaśāstraṁ}
  \item \emph{Recall:} say the Telugu for history, then the Telugu for geography, then say \emph{vijñānaśāstraṁ} again
\end{itemize}

\begin{tcolorbox}[breakable,colback=teal!4,colframe=teal!35!black,title={Before you move on}]
What does \te{విజ్ఞానశాస్త్రం} mean? (\textquotedblleft{}Science\textquotedblright{}.) Say it once more.
\end{tcolorbox}
\section[\texorpdfstring{\te{భౌతికశాస్త్రం} (bhautikaśāstraṁ)}{భౌతికశాస్త్రం (bhautikaśāstraṁ)}]{\te{భౌతికశాస్త్రం} (bhautikaśāstraṁ) — physics}

\label{lesson:TE-C238-bhautikasastram}



\begin{quote}
Before the new one: say the Telugu for geography, then the Telugu for science.
\end{quote}

\subsection*{You'll want to know: \te{భౌతికశాస్త్రం}}
\textbf{\te{భౌతికశాస్త్రం}} — \emph{bhautikaśāstraṁ} — \textquotedblleft{}physics\textquotedblright{}.

\begin{grammarlens}[title={What you've built}]
One more word for this chapter.
\end{grammarlens}

\subsection*{Your turn}
\begin{itemize}
\raggedright
  \item \emph{Say it:} \emph{bhautikaśāstraṁ}
  \item \emph{Say it:} \emph{bhautikaśāstraṁ}, once more
  \item \emph{Say it:} say \emph{bhautikaśāstraṁ}
  \item \emph{Recall:} say the Telugu for geography, then the Telugu for science, then say \emph{bhautikaśāstraṁ} again
\end{itemize}

\begin{tcolorbox}[breakable,colback=teal!4,colframe=teal!35!black,title={Before you move on}]
What does \te{భౌతికశాస్త్రం} mean? (\textquotedblleft{}Physics\textquotedblright{}.) Say it once more.
\end{tcolorbox}
